\documentclass[11pt,UTF8,fontset=fandol]{ctexart}
\usepackage{amsmath,amssymb,amsfonts}
\usepackage{graphicx,float,caption,subcaption}
\usepackage{booktabs,array,multirow,makecell}
\usepackage{geometry,enumitem,xcolor}
\usepackage[hidelinks]{hyperref}
\usepackage{url}
\usepackage{tcolorbox}
\geometry{margin=2.3cm}
\graphicspath{{/Users/yye/project/论文/assets/ragen/}}
\captionsetup{font=small,labelfont=bf}
\setlength{\parskip}{0.45em}

% 知识卡片小盒子
\newtcolorbox{cardbox}[1][]{colback=blue!4,colframe=blue!40!black,
  boxrule=0.5pt,arc=2pt,left=6pt,right=6pt,top=4pt,bottom=4pt,#1}
% 要点/发现盒子
\newtcolorbox{findbox}[1][]{colback=orange!6,colframe=orange!55!black,
  boxrule=0.5pt,arc=2pt,left=6pt,right=6pt,top=4pt,bottom=4pt,#1}

\title{\textbf{RAGEN：从多轮强化学习视角理解 LLM 智能体的自我进化\\——一份面向 AI 应用开发者的中文精读导读}}
\author{Zihan Wang, Kangrui Wang, Qineng Wang, Pingyue Zhang, Linjie Li, Zhengyuan Yang 等\\
\small Northwestern University（西北大学）\,/\,University of Washington（华盛顿大学）\,/\,Stanford University（斯坦福大学）\\
\small Microsoft（微软）\,/\,New York University（纽约大学）\,/\,University of British Columbia\,/\,Singapore Management University}
\date{}

\begin{document}
\maketitle

\begin{center}
\fbox{\parbox{0.95\linewidth}{\small
\textbf{英文原题}：RAGEN: Understanding Self-Evolution in LLM Agents via Multi-Turn Reinforcement Learning\\
\textbf{发表}：arXiv 2025（arXiv:2504.20073）\quad\textbf{引用数}：229（Semantic Scholar，2026-06-07）\\
\textbf{代码/环境}：\url{https://github.com/RAGEN-AI/RAGEN}\quad\textbf{项目主页}：\url{https://ragen-ai.github.io/}\\[2pt]
\textbf{一句话定位}：这\emph{不是}又一个刷榜的新方法，而是认真研究"多轮 agentic RL（智能体强化学习）为什么训着训着就崩"，并给出一套\textbf{稳定训练配方}。作者刻意在极简、可控的玩具环境里训练，把训练动力学（training dynamics）看个清楚。\textbf{想自己动手训 agent 的人，避坑必读。}
}}
\end{center}

\begin{center}
\small\textit{【阅读说明】本文是面向 AI 应用开发实习生/初学者的"精读导读"，属于转化性学习材料：以讲解、转述、关键句翻译为主，并非对原文的逐字翻译。术语首次出现时括注英文原文；模型名、数据集名、公式、引用标记按规范保留原样。}
\end{center}

%==================================================================
\section{小白知识卡：读这篇前要懂的背景}
%==================================================================

这篇论文站在"强化学习 + LLM 智能体"的交叉口上，下面这些概念先有个大白话印象，后面读起来会顺很多。

\begin{cardbox}
\textbf{强化学习（RL，reinforcement learning）}：让模型像玩游戏一样，通过"试一下、看奖励、再调整"来学习。核心三要素是\textbf{状态（state）}、\textbf{动作（action）}、\textbf{奖励（reward）}。
\end{cardbox}

\begin{cardbox}
\textbf{单轮 RL vs.\ 多轮 RL（single-turn vs.\ multi-turn）}：单轮就是"给一道数学题、答一次、对错给分"，一问一答就结束（如解数学题、写代码）。多轮是"模型要连续做好几步决策，每步环境都会反馈、状态会变"（如推箱子、网购、操作 GUI）。本文的核心就是研究\textbf{多轮}场景。
\end{cardbox}

\begin{cardbox}
\textbf{LLM 智能体（LLM agent）}：把大语言模型当成一个能在环境里\textbf{自主行动}的"代理人"，它要做序列决策、跨轮记忆、并适应带随机性的环境反馈。
\end{cardbox}

\begin{cardbox}
\textbf{PPO（Proximal Policy Optimization，近端策略优化）}：最经典的 RL 算法之一。它额外训练一个\textbf{critic（评论家/价值网络）}来估计"当前局面值多少分"，并用"裁剪（clipping）"限制每步更新幅度，让训练更稳。
\end{cardbox}

\begin{cardbox}
\textbf{GRPO（Group Relative Policy Optimization，组相对策略优化）}：DeepSeek 提出的算法，\textbf{不用 critic}。它对同一道题采样一组答案，用"这组答案里你比平均水平好多少"来当优势（advantage）。更省，但少了 critic 的稳定作用。
\end{cardbox}

\begin{cardbox}
\textbf{RLVR（基于可验证奖励的强化学习，RL with Verifiable Reward）/ 规则奖励（rule-based reward）}：奖励不靠人打分、也不靠奖励模型，而是用\textbf{规则}自动判定对错（如数学答案是否正确、箱子是否推到位）。本文用的就是这种规则奖励。
\end{cardbox}

\begin{cardbox}
\textbf{rollout（轨迹采样/推演）}：让当前模型在环境里实际跑一遍，生成一条完整的"状态-动作-奖励"轨迹（trajectory）。这些自己跑出来的轨迹，就是 RL 训练的"教材"。
\end{cardbox}

\begin{cardbox}
\textbf{credit assignment（信用分配）}：一条多轮轨迹最后成功/失败了，到底是中间\emph{哪一步}的功劳或锅？把最终奖励合理地"分摊"到每一步上，是多轮 RL 的核心难题。
\end{cardbox}

\begin{cardbox}
\textbf{熵（entropy）与奖励方差（reward variance）}：熵衡量模型输出有多"随机/多样"；奖励方差衡量同一道题的多次尝试结果有多"参差"。两者都高，说明模型还在积极探索；两者一塌缩，往往就是"学僵了"的信号——这正是本文诊断训练崩溃的关键仪表盘。
\end{cardbox}

%==================================================================
\section{摘要全译}
%==================================================================

\begin{quote}\small
把大语言模型（LLM）训练成交互式智能体面临独特挑战，包括长时程决策（long-horizon decision making）以及与带随机性的环境反馈进行交互。尽管强化学习（RL）已在静态任务上取得进展，多轮智能体 RL 训练仍未被充分探索。我们提出 \textbf{StarPO}（State-Thinking-Actions-Reward Policy Optimization，状态-思考-动作-奖励策略优化），一个面向轨迹级（trajectory-level）智能体 RL 的通用框架；并引入 \textbf{RAGEN}，一个用于训练和评估 LLM 智能体的模块化系统。我们在四个风格化（stylized）环境上的研究揭示了三个核心发现。第一，我们的智能体 RL 训练表现出一种反复出现的失效模式——\textbf{Echo Trap（回声陷阱）}，其特征是奖励变化骤降（cliff）与梯度尖峰（gradient spike）；我们用 \textbf{StarPO-S}（稳定版变体）来应对它，包含轨迹过滤（trajectory filtering）、引入 critic、以及梯度稳定化。第二，我们发现 RL rollout 的塑造会从以下方面受益：多样的初始状态、适中的交互粒度（interaction granularity）、以及更频繁的采样。第三，我们表明，如果没有细粒度、感知推理（reasoning-aware）的奖励信号，智能体的推理能力很难通过多轮 RL 涌现出来，模型可能表现出肤浅的策略或幻觉式的思考（hallucinated thoughts）。

\textbf{关键词}：LLM Agents，Multi-turn RL\quad\textbf{主页}：\url{https://ragen-ai.github.io/}\quad\textbf{代码/环境}：\url{https://github.com/RAGEN-AI/RAGEN}
\end{quote}

%==================================================================
\section{逐节精读}
%==================================================================

\subsection{第 1 节：引言——核心问题是什么}

\textbf{一句话讲清}：把 LLM 训成能在交互环境里"自我进化"的智能体，最大的拦路虎是\textbf{训练不稳定}。这篇就是系统研究"到底什么设计因素能让自我进化的 LLM 智能体学得又好又稳"。

作者先点出多轮 agent 和静态任务（如单轮数学题、写代码）的本质差别：agent 要\textbf{连续做决策、跨轮维持记忆、适应随机反馈}。这类场景（规划助手、机器人、辅导 agent）不仅要求模型"表现好"，更要求它\textbf{能从经验中自我提升}。

原文一句关键设问值得记住：
\begin{quote}\small
"A key open question is: what design factors make self-evolving LLM agents learn effectively and stably?"\\
（一个关键的开放问题是：什么样的设计因素能让自我进化的 LLM 智能体学得既有效又稳定？）
\end{quote}

为回答这个问题，作者搭了两样东西：\textbf{StarPO}（一个把"多轮、轨迹级"训练统一起来的 RL 框架，可灵活控制推理、奖励分配、prompt-rollout 结构）和 \textbf{RAGEN}（在 StarPO 之上实现完整训练闭环的模块化系统，专门用来做训练动力学的系统性分析）。

评测放在四个复杂度递增的环境上：\textbf{Bandit}（单轮、随机）、\textbf{Sokoban}（多轮、确定性）、\textbf{Frozen Lake}（多轮、随机）、\textbf{WebShop}（多轮、开放域）。前三个是极简、完全可控的符号环境，方便"看清"动力学；WebShop 加入了真实世界的语言理解和推理。

引言最后用三条"发现"概括全文（这也是整篇的骨架，后面会逐一展开）：

\begin{findbox}
\textbf{发现 1（梯度稳定是稳定训练的关键）}：多轮 RL 常陷入一种反复出现的不稳定模式 \textbf{Echo Trap（回声陷阱）}——智能体过拟合到\emph{局部被奖励}的推理套路，表现为奖励方差塌缩、熵下降、梯度尖峰。对策是 \textbf{StarPO-S}：基于方差的轨迹过滤、critic 基线、以及解耦裁剪（decoupled clipping）。
\end{findbox}

\begin{findbox}
\textbf{发现 2（rollout 的频率与多样性塑造自我进化）}：LLM 自己生成的轨迹是核心训练材料。关键的 rollout 因素有三个：（1）rollout 要来自\textbf{多样的初始状态}，且每个初始状态有多条回应；（2）每轮允许\textbf{多个动作}，在固定轮数内拉长交互时程；（3）保持\textbf{高 rollout 频率}，让在线反馈反映当前策略。
\end{findbox}

\begin{findbox}
\textbf{发现 3（推理的涌现需要精心设计的奖励信号）}：仅仅在输出格式上"鼓励推理"并不能保证模型真的去推理。即便用 \texttt{<think>} 标签提示、并用 StarPO 做轨迹级优化，如果推理\emph{带不来额外奖励优势}，模型就会退化回"直接选动作"。而当奖励只反映任务是否成功时，模型会产生\textbf{幻觉式推理}——想的和环境状态对不上。这说明长时程 agent 训练需要\textbf{细粒度、感知推理}的奖励设计。
\end{findbox}

\subsection{第 2 节：框架——StarPO 是怎么定义的}

\subsubsection{2.1 把 agent 训练写成 MDP}

以前给语言模型做 RL 通常是\textbf{单轮}设定：在数据集 $\mathcal{D}$ 上采样 prompt-response 对 $(s,a)$，目标是最大化期望奖励：
\begin{equation*}
J_{\text{step}}(\theta)=\mathbb{E}_{s\sim\mathcal{D},\,a\sim\pi_\theta(\cdot\,|\,s)}\big[R(s,a)\big].\tag{1}
\end{equation*}

\textbf{逐符号拆解（式 1）}：$\theta$ 是模型参数；$\pi_\theta$ 是\textbf{策略}（policy，就是当前模型）；$s$ 是输入 prompt（状态）；$a$ 是模型采样出的回应（动作）；$R(s,a)$ 是奖励函数；$\mathbb{E}[\cdot]$ 表示"求平均/期望"。整条式子的意思是：\emph{让模型在所有 prompt 上、按自己的策略采样回应，平均奖励越高越好}。

但 LLM 智能体得在\textbf{多轮、带随机反馈}的环境里干活。作者把它建模成马尔可夫决策过程（MDP，Markov Decision Process）$\mathcal{M}=\{S,A,P\}$：$S$ 是状态（如观测序列或交互历史），$A$ 是动作（通常是 token 序列），$P$ 是转移动力学与奖励生成过程。在每个时间步 $t$，策略基于当前状态和历史生成动作，环境再返回奖励和新状态：
\begin{equation*}
a_t\sim\pi_\theta(\cdot\,|\,s_t,\tau_{<t}),\qquad (r_t,s_{t+1})\sim P(\cdot\,|\,s_t,a_t).
\end{equation*}
其中 $\tau_{<t}=\{s_0,a_0,r_0,\dots,s_{t-1},a_{t-1},r_{t-1}\}$ 是\textbf{交互历史}。这个过程持续到最大时程 $K$，产出一条完整轨迹 $\tau=\{s_0,a_0,r_0,\dots,s_K\}$，它就是智能体学习的"教材"。

\subsubsection{2.2 StarPO：用轨迹级优化来强化推理}

\textbf{核心思想（最重要的一段）}：和以前"把每个动作独立对待"的静态任务方法不同，StarPO 把\textbf{整条轨迹}——包括观测、推理痕迹、动作、反馈——当作一个\textbf{连贯整体}来 rollout 和优化。目标是最大化期望\emph{轨迹}奖励：
\begin{equation*}
J_{\text{StarPO}}(\theta)=\mathbb{E}_{\mathcal{M},\,\tau\sim\pi_\theta}\big[R(\tau)\big].\tag{2}
\end{equation*}

\textbf{逐符号拆解（式 2）}：$\mathcal{M}$ 是那个 MDP；$\tau$ 是一条"被推理增强过"的完整交互序列；$R(\tau)$ 是\textbf{整条轨迹}的累积奖励。轨迹概率 $\pi_\theta(\tau)$ 可以分解成 token 级的似然，所以 StarPO 天然兼容自回归 LLM。

\textbf{和老方法的对比一眼看懂}：

\begin{tcolorbox}[colback=gray!5,colframe=gray!50,boxrule=0.5pt,arc=2pt]
\small
\textbf{老方法}（PPO、GRPO）：$J_{\text{step}}(\theta)=\mathbb{E}_{x\sim\mathcal{D},\,y\sim\pi_\theta(\cdot|x)}[R(x,y)]$ —— 给输入 $x$，优化\textbf{单轮输出} $y$。\\[3pt]
\textbf{StarPO（本文）}：$J_{\text{StarPO}}(\theta)=\mathbb{E}_{\mathcal{M},\,\tau\sim\pi_\theta}[R(\tau)]$ —— 优化\textbf{整条轨迹} $\tau=\{s_0,a_0,r_0,\dots,s_K\}$ 上的总奖励。
\end{tcolorbox}

\textbf{优化流程}：每轮训练，智能体从初始状态 $s_0$ 出发，生成 $N$ 条轨迹。每一步产出一个\textbf{被推理引导的结构化输出}：
\begin{equation*}
a^T_t=\texttt{<think>}\dots\texttt{</think>}\texttt{<answer>}\,a_t\,\texttt{</answer>},\tag{3}
\end{equation*}
这里 $a^T_t$ 是包含中间推理的完整动作输出，而 $a_t$ 是真正能在环境里执行的子动作序列。环境随后返回下一状态 $s_{t+1}$ 和奖励 $r_t$。整条轨迹里每个部分要么是 LLM 生成的、要么是环境产生的，都会被\textbf{联合优化}。

训练循环交替进行 rollout 和 update：新轨迹可用 $\pi_\theta$ 在线（on-policy）生成，也可从 replay buffer 里按 $\pi_{\text{old}}$ 采样。每个训练循环有 $P$ 个初始状态、每个生成 $N$ 条轨迹、以批大小 $E$ 更新、共 $L$ 个循环，总梯度更新步数 $S=L\cdot\frac{P\cdot N}{E}$。

\subsubsection{2.2.3 模块化的优化策略}

StarPO 在统一的轨迹级抽象下支持多种优化算法。论文给了两种 token 级更新的实例：

\textbf{PPO 版本}（训练一个 critic 来估计 token 级价值和优势 $A_{i,t}$）：
\begin{equation*}
J_{\text{PPO}}(\theta)=\frac{1}{G}\sum_{i=1}^{G}\frac{1}{|\tau_i|}\sum_{t=1}^{|\tau_i|}\min\!\Big(\tfrac{\pi_\theta(\tau_{i,(t)}|\tau_{i,<t})}{\pi_{\text{old}}(\tau_{i,(t)}|\tau_{i,<t})}A_{i,t},\ \text{clip}\big(\tfrac{\pi_\theta(\tau_{i,(t)}|\tau_{i,<t})}{\pi_{\text{old}}(\tau_{i,(t)}|\tau_{i,<t})},1-\varepsilon,1+\varepsilon\big)A_{i,t}\Big).\tag{4}
\end{equation*}

\textbf{GRPO 版本}（不用 critic，给每条轨迹一个标量奖励 $R(\tau_i)$，并在轨迹内所有 token 上做归一化优势 $\hat{A}_{i,t}$）：
\begin{equation*}
\hat{A}_{i,t}=\frac{R(\tau_i)-\text{mean}(\{R(\tau_1),\dots,R(\tau_G)\})}{\text{std}(\{R(\tau_1),\dots,R(\tau_G)\})}.\tag{5}
\end{equation*}
再代入与式(4)结构相同的 GRPO 目标 $J_{\text{GRPO}}(\theta)$（把 $A_{i,t}$ 换成 $\hat{A}_{i,t}$，见式 6）。

\textbf{怎么读这几个公式（直觉版）}：$G$ 是这批里的轨迹数；$\tau_{i,(t)}$ 是第 $i$ 条轨迹的第 $t$ 个 token；那个分式 $\frac{\pi_\theta}{\pi_{\text{old}}}$ 是\textbf{新旧策略的概率比}（这一步比之前更愿意/更不愿意产出这个 token）；$A_{i,t}$（或 $\hat A_{i,t}$）是\textbf{优势}（这一步比平均水平好多少，好就拉高概率、差就压低）；$\text{clip}(\cdot,1-\varepsilon,1+\varepsilon)$ 把概率比裁剪在 $[1-\varepsilon,1+\varepsilon]$ 内，\textbf{防止一步更新太猛}。两者的根本区别就在优势怎么算：PPO 靠 critic 估计，GRPO 靠"组内减均值除标准差"。

\subsubsection{2.3 RAGEN 系统}

RAGEN 是 StarPO 的工程落地：支持结构化 rollout、可定制奖励函数、与多轮随机环境的对接。它既是 StarPO 的执行后端，也是研究稳定性、泛化、学习动力学的平台。设计上强调\textbf{可扩展}——新环境、新奖励方案、新 rollout 策略都能即插即用。

\subsection{第 3 节：实验设置}

\textbf{四个环境，各考一种能力}（细节见后文"实验解读"）：Bandit 考"噪声反馈下的风险敏感推理"；Sokoban 考"不可逆的符号规划"；Frozen Lake 把规划和概率性转移结合；WebShop 涉及自然语言落地和网页交互。前三个符号环境刻意做得\textbf{极简、完全可控}，就是为了"看清"训练过程。

\textbf{训练设置}：主实验用 Qwen-2.5 Instruct 0.5B 训三个符号任务、用 3B 变体训更难的 WebShop。在 H100 GPU 上跑 100--200 个 rollout--update 迭代。每批采样 $P{=}8$ 个 prompt，每个 prompt 出 $N{=}16$ 条 rollout，最多 5 轮、10 个动作。策略更新用 GRPO 或 PPO + GAE（$\gamma{=}1.0,\lambda{=}1.0$）、Adam 优化器、熵奖励（$\beta{=}0.001$）、以及一个 $-0.1$ 的格式惩罚（response-format penalty）。

\textbf{评测指标（五个仪表盘）}：每个环境用 256 个固定 prompt、温度 $T{=}0.5$、5 轮后截断。指标包括：(i) 成功率（任务完成）；(ii) rollout 熵（探索程度）；(iii) 组内奖励方差（行为多样性）；(iv) 回应长度（推理冗长度）；(v) 梯度范数（训练稳定性）。\textbf{后四个不是为了刷分，而是诊断训练健康度的"体检仪表"}——这正是这篇方法论文的精髓。

%==================================================================
\section{方法图解：StarPO 框架到底在干什么}
%==================================================================

\begin{figure}[H]\centering
  \includegraphics[width=\linewidth]{fig_starpo.pdf}
  \caption{StarPO（State-Thinking-Actions-Reward Policy Optimization）框架（原文图 2）。LLM 生成"推理引导的动作"与环境进行多轮交互，累积轨迹级奖励，归一化后用于更新 LLM 策略。}
  \label{fig:starpo}
\end{figure}

\textbf{这张图在说什么、数据怎么流动}（对照图~\ref{fig:starpo}，分三块看）：

\textbf{左半部分 —— K 轮 Rollout（K-turn Rollout）}：这是"模型在环境里实际跑一遍"的过程。
\begin{itemize}[leftmargin=1.4em,itemsep=2pt,topsep=2pt]
\item 从初始状态 $s_0$ 开始，\textbf{LLM-Rollout} 模块 \emph{generate（生成）} 一个结构化输出（中间那条）：\texttt{<think>} 思考过程 \texttt{</think>}\,\texttt{<ans>} 真正动作 $a_K$ \texttt{</ans>}。
\item 这个动作序列（action sequence）交给 \textbf{Env（环境，图里那个小房子）} 执行，环境 \emph{feedback（反馈）} 回奖励 $r_K$ 和新状态 $s_{K+1}$。
\item 新的 $(a^T_K, r_K, s_{K+1})$ 被 \emph{append（追加）} 回输入序列，进入下一轮。如此循环 $K$ 轮。
\item 最终把所有片段拼成一条完整\textbf{轨迹 $\tau$}（最下面那条带子）：$s_0,a^T_0,r_0,s_1,\dots$，并算出整条轨迹的总奖励 $r_{\text{all}}=\sum_i r_i$。
\end{itemize}

\textbf{右半部分 —— 更新 LLM（Update LLM）}：这是"从跑出来的轨迹里学习"的过程。
\begin{itemize}[leftmargin=1.4em,itemsep=2pt,topsep=2pt]
\item 从同一个 $s_0$ 出发，会有 $N$ 条平行轨迹（Trajectory 1、2、$\dots$、$N$），每条都有自己的奖励 $R^{(i)}$。
\item 这 $N$ 条轨迹一起送进 \textbf{Advantage Estimation（优势估计）} 模块，算出每条轨迹的优势 $A^{(i)}$——本质就是"这条比这一组的平均水平好多少"。
\item 最后用这些优势去做 \textbf{LLM-Update}，优化目标 $\max_\theta J_{\text{StarPO}}(\theta,R)$，把"表现好的轨迹"对应的 token 概率拉高。
\end{itemize}

\textbf{一句话总结这张图}：\emph{左边"跑"出多条完整轨迹、右边按"组内相对好坏"给轨迹打优势分、再回去更新模型}——这就是"轨迹级"优化区别于"单步"优化的全部精髓。

%==================================================================
\section{实验解读：数字背后说明了什么}
%==================================================================

\subsection{四个测试环境（数据集/基准）}

\begin{itemize}[leftmargin=1.4em,itemsep=2pt,topsep=2pt]
\item \textbf{Bi-Arm Bandit（双臂老虎机，单轮、随机）}：在两个有"符号语义"的选项间二选一（如 Dragon vs.\ Phoenix）。低风险臂固定给 0.15 的奖励；高风险臂从 $\text{Bernoulli}(0.25)$ 采样——\textbf{方差更大但期望更高}。妙处在于：低风险臂\emph{单次}更常赢，但高风险臂\emph{长期}更优。这就逼着模型去做风险敏感推理，把符号线索（"Dragon"）和背后的奖励统计关联起来，而不是被短期信号带偏。
\item \textbf{Sokoban（推箱子，多轮、确定性）}：在格子里把箱子推到目标点。关键是\textbf{不可逆}——箱子只能推不能拉，推错就死局，逼模型提前规划。奖励：每个到位的箱子 $+1$，推偏 $-1$，完成 $+10$，每个动作 $-0.1$。
\item \textbf{Frozen Lake（冰湖，多轮、随机）}：在有"滑冰格"的格子里走到目标。每个动作只有 $1/3$ 概率成功，$2/3$ 概率向垂直方向滑偏。奖励稀疏：成功 $+1$，其余全 $0$。
\item \textbf{WebShop（网购，多轮、开放域）}：真实的多轮网购任务，要理解自然语言需求、操作半结构化界面、搜索/点击/读商品描述。引入了符号环境里没有的语言落地和动作空间挑战。
\end{itemize}

\textbf{对比的基线}：本文的"基线"不是别家的系统，而是把\textbf{单轮 RL 方法（PPO、GRPO）直接搬到多轮}（即 vanilla StarPO），再和稳定版 \textbf{StarPO-S} 对比；同时附录里还和更大模型（72B）、前沿模型（GPT-4o、Qwen-2.5-72B）对照，以便给小模型的表现一个参照系。

\subsection{核心结果一：直接搬单轮方法会"训崩"（发现 1\&2\&3）}

baseline StarPO 在符号任务（Bandit、Sokoban）上\textbf{先涨后崩}：早期有提升，最终塌缩。PPO 比 GRPO 稳一些（崩得更晚、峰值更高），因为它的 critic 提供更平滑的奖励估计；但在 Frozen Lake 上反而 GRPO 更稳（因为这个任务状态价值难估，会把 PPO 带乱）；WebShop 上两者都成功（强语言先验 + 高初始奖励，降低了对 critic 的依赖）。

作者对比"早期 vs.\ 晚期"轨迹，发现了 \textbf{Echo Trap（回声陷阱）}：早期模型对符号意义和期望奖励有\textbf{多样}的推理，晚期回应变得\textbf{重复、确定}。这说明 RL 把"局部被奖励的套路"过度放大了，反复套用记住的模板、压制了探索。原文一句关键描述：
\begin{quote}\small
"early-stage agent respond with diverse symbolic reasoning, but collapse into deterministic, repetitive templates after training."\\
（早期智能体会用多样的符号推理来回应，但训练后塌缩成确定、重复的模板。）
\end{quote}

\textbf{怎么提前预警崩溃}（这是全文最实用的部分之一）：作者监控四个信号，并发现它们有明确的时间先后：
\begin{itemize}[leftmargin=1.4em,itemsep=2pt,topsep=2pt]
\item \textbf{奖励标准差（reward std）是收敛的早期指标}：它常在奖励均值崩之前就先掉。例如 FrozenLake-PPO 在第 40 步 std 骤降，而奖励均值第 90 步才崩。
\item \textbf{梯度范数尖峰（gradient norm spike）标志不可逆崩溃}：一旦出现尖峰（如 Bandit 第 170 步、Sokoban 第 110 步、FrozenLake 第 90 步），哪怕极小的参数更新也会引发剧烈损失波动，之后基本救不回来。
\item \textbf{熵（entropy）在有效学习时应平稳衰减}：熵的急升或乱跳，往往对应推理行为已经崩坏。
\end{itemize}

\begin{findbox}
\textbf{发现 1\&2\&3 合起来的实操结论}：单轮 RL 方法不能直接拿来做多轮 agent RL；模型崩溃在训练中表现为 Echo Trap；而\textbf{奖励标准差/熵的异常波动}会先于性能下降出现，\textbf{梯度范数尖峰}则标志着不可逆的崩溃点——这几个指标给了我们"在崩之前踩刹车"的早期信号。
\end{findbox}

\subsection{核心结果二：StarPO-S 怎么把训练救回来（发现 4）}

基于"奖励标准差下降常常预示崩溃"这个洞见，作者提出一个反直觉但有效的假设：\emph{是不是应该把训练资源更多投在那些"行为更不确定、奖励方差更高"的任务实例上？}——这正是\textbf{主动学习（Active Learning）}的思想：不确定的样本信息量最大。

为此定义\textbf{轨迹级结果不确定性 $U$}（对初始状态 $s_0$）：
\begin{equation*}
U(\pi_\theta,\mathcal{M},s_0)=\text{Std}_{\tau\sim\pi_\theta(\cdot\,|\,s_0)}\big[R(\tau)\big].\tag{7}
\end{equation*}
\textbf{拆解（式 7）}：对同一个初始状态 $s_0$ 反复 rollout 多条轨迹，看它们最终奖励 $R(\tau)$ 的\textbf{标准差}。std 大 = 这道题模型时对时错、最值得学；std 小 = 要么稳过、要么稳挂，信息量低。训练时按这个 std 给 prompt 排序，\textbf{只保留方差最高的 top $p\%$}。

\textbf{StarPO-S 的三件套}：
\begin{enumerate}[leftmargin=1.6em,itemsep=2pt,topsep=2pt]
\item \textbf{基于不确定性的轨迹过滤}：扔掉低方差 rollout。PPO 上效果显著——保留 75\% 能把 FrozenLake 的稳定期从 100 步延到 140 步，保留 50\% 干脆完全不崩。默认采用激进的 \textbf{保留 25\%}（但作者也提醒这未必对所有场景最优）。过滤还顺带\textbf{提速}（少算了没信息量的样本）。
\item \textbf{KL 项移除（KL Term Removal）}：借鉴 DAPO，去掉 KL 惩罚项，给优化更大自由度。
\item \textbf{Clip-Higher（非对称裁剪，Asymmetric Clipping）}：也来自 DAPO，把裁剪上下界解耦，鼓励对低概率但有潜力的 token 多一点探索。
\end{enumerate}

\begin{findbox}
\textbf{发现 4}：在\textbf{高方差 prompt} 上训练，能延迟甚至消除多轮 RL 的崩溃。StarPO-S 通过丢弃低信息量 rollout，既\textbf{提升性能}又\textbf{减少更新步数}（PPO 上尤其明显），与主动学习原理一致。整体对比（原文图 6）显示，StarPO-S 在全部四个任务上都\textbf{延迟了崩溃、并取得更高成功率}。
\end{findbox}

\subsection{核心结果三：什么样的 rollout 才是"好教材"（发现 5）}

作者在 vanilla Sokoban 上训练、在变体任务（SokobanNewVocab、LargeSokoban、FrozenLake）上评测，研究三个 rollout 维度。下面把关键数字重排成小表。

\begin{table}[H]\centering
\caption{任务多样性对泛化的影响（\%）（原文表 1）。固定批大小下，"每个 prompt 的回应数"越少 = 不同 prompt 越多 = 多样性越高。结论：每个 prompt 出 \textbf{4} 条回应时泛化最好。}
\begin{tabular}{cccc}
\toprule
每个 prompt 的回应数 & SingleSokoban & SokobanNewVocab & FrozenLake \\ \midrule
32 & 21.09 & 20.22 & 17.97 \\
16 & 20.31 & 21.48 & 19.53 \\
8  & 20.31 & 19.53 & 17.19 \\
\textbf{4}  & \textbf{20.70} & \textbf{25.39} & \textbf{21.48} \\
2  & 19.92 & 25.00 & 12.50 \\
1  & 19.53 & 22.27 & 12.50 \\
\bottomrule
\end{tabular}
\end{table}

\textbf{解读表 1}：多样性高（不同初始状态多）泛化更好，\textbf{但前提是每个 prompt 仍要有多条 rollout}（否则没法"对比不同结果")。所以"4 条/prompt"是甜区：既保证了 prompt 多样，又保留了组内对比；而"1 条/prompt"两端都垮（FrozenLake 只有 12.50\%）。

\begin{table}[H]\centering
\caption{不同"每轮动作预算"下各环境的表现（\%）（原文表 2）。\textbf{5--6} 个动作/轮最好，给规划留足空间又不至于太乱。}
\begin{tabular}{ccccc}
\toprule
每轮最大动作数 & Sokoban & SokobanNewVocab & LargeSokoban & FrozenLake \\ \midrule
1 & 12.11 & 13.67 & 1.17 & 11.72 \\
2 & 16.41 & 21.09 & 3.52 & 18.36 \\
3 & 19.53 & 19.53 & 1.95 & 20.88 \\
4 & 26.95 & 26.95 & 5.08 & 20.70 \\
5 & 28.13 & 25.78 & 6.25 & 21.09 \\
\textbf{6} & \textbf{33.59} & \textbf{31.64} & \textbf{6.64} & 18.36 \\
7 & 22.27 & 28.52 & 3.91 & 19.53 \\
\bottomrule
\end{tabular}
\end{table}

\textbf{解读表 2}：动作预算太小（1--2）规划空间不够；太大（7）则注入噪声、稀释奖励反馈。\textbf{适中（5--6）}最佳——这是一个明确的"交互粒度"调参建议。

\textbf{第三个维度——rollout 频率（Online-$k$）}：用一组 rollout 连续更新 $k$ 次策略，$k$ 越小越"新鲜"。Online-1 即完全在线（每次更新都用新鲜 rollout）。结果（原文图 7）：\textbf{越新鲜（Online-1）收敛越快、泛化越好}；用陈旧数据（Online-5/10）则变差。

\begin{findbox}
\textbf{发现 5}：任务多样性、动作预算、rollout 频率共同决定 rollout 质量。\textbf{多样的初始状态}带来更好的策略对比和泛化；\textbf{适中的动作预算}给足规划空间又不引入噪声；\textbf{够新鲜的 rollout}让优化目标和当前策略保持对齐。
\end{findbox}

\subsection{核心结果四：推理能帮泛化，但在多轮里会"自己消失"（发现 6）}

这是全文最发人深省的一节。作者设计了一个受控泛化测试：在 Bandit 上用 [Teacher, Engineer] 这对臂训练、在 [Librarian, Trader] 上评测（保持"高风险高回报"的直觉对应）；再做一个 \textbf{BanditRev}（把符号-奖励关系\textbf{反转}，让推理变得反直觉、更难）。

\begin{table}[H]\centering
\caption{有/无推理时的泛化表现（\%），均在 StarPO-S 下（原文表 3）。关掉推理（NoThink）在\textbf{单轮 Bandit} 上显著掉点，但在\textbf{多轮 Sokoban} 上影响很小甚至混杂。}
\begin{tabular}{lccccc}
\toprule
 & \multicolumn{3}{c}{在 Bandit 上训练} & \multicolumn{2}{c}{在 Sokoban 上训练} \\
\cmidrule(lr){2-4}\cmidrule(lr){5-6}
 & Bandit & Bandit-Rev & FrozenLake & LargeSokoban / Sokoban & SokobanNewVocab \\ \midrule
StarPO-S & \textbf{100.00} & \textbf{67.58} & 19.92 & 2.34 / 21.48 & 18.75 \\
NoThink  & 81.25 & 56.25 & 19.53 & 2.73 / 20.73 & \textbf{26.17} \\
\bottomrule
\end{tabular}
\end{table}

\textbf{解读表 3}：在\textbf{单轮} Bandit，带推理痕迹的模型泛化明显更好（100\% vs.\ 81.25\%），即便在反直觉的 BanditRev 也是（67.58\% vs.\ 56.25\%）——说明推理帮模型\textbf{内化了符号-奖励关联}，而非死记硬背。但在\textbf{多轮} Sokoban，去掉推理（NoThink）反而打平甚至更好（SokobanNewVocab 上 26.17\% \textgreater\ 18.75\%）。

为什么多轮里推理会失效？看推理长度的变化：

\begin{table}[H]\centering
\caption{不同训练步数下的推理长度（\texttt{<think>} 块的 token 数）（原文表 4）。token 长度普遍随训练\textbf{下降}，而"上下文更难"的 ReverseBandit 需要的推理更多、衰减更慢。}
\begin{tabular}{lcccccc}
\toprule
 & \multicolumn{4}{c}{Bandit} & \multicolumn{2}{c}{Sokoban} \\
\cmidrule(lr){2-5}\cmidrule(lr){6-7}
Step & Original & NoThink & Reverse & Reverse-NoThink & Original & NoThink \\ \midrule
0   & 66.0 & 12.7 & 68.8 & 12.7 & 307.1 & 68.6 \\
100 & 25.6 & 12.4 & 33.7 & 13.0 & 104.6 & 55.3 \\
200 & 17.6 & 12.4 & 30.7 & 13.0 & 89.5  & 60.0 \\
\bottomrule
\end{tabular}
\end{table}

\textbf{解读表 4}：推理痕迹随训练\textbf{持续缩短}（如 Sokoban Original 从 307 token 掉到 89），说明\textbf{模型在主动抑制自己的思考过程}。有意思的是，在"语义错位、必须靠推理"的 BanditRev 里，推理保持得更长——印证了"上下文越难、推理越被保留"。

\textbf{根因（作者的推断）}：多轮任务的奖励\textbf{稀疏、延迟}，常常分不清"是靠连贯推理成功的"还是"瞎试碰对的"。附录里就有例子：模型生成了\textbf{不连贯/幻觉式的推理，却照样拿高奖励}。这就引出一个核心难题：\emph{当奖励本身无法反映推理质量时，怎么才能稳定地强化"有用的推理"？}

\begin{findbox}
\textbf{发现 6}：符号推理能在简单的单轮任务里、弱监督下涌现，但在多轮环境里\textbf{若奖励不显式鼓励可解释的中间推理步骤，推理就会逐渐衰减}。一个初步对策是用\textbf{基于格式的惩罚}（输出缺少有效的 \texttt{<think>--<answer>} 结构就扣分）；作者呼吁未来探索更细粒度的奖励设计（如奖励"部分正确"）来可靠地强化长时程决策中的推理。
\end{findbox}

\subsection{第 5\&6 节：相关工作与结论/局限}

\textbf{相关工作}一句话带过：本文站在两条线的交叉口——一条是"用 RL 提升 LLM 推理"（PPO、GRPO、DAPO、STaR、PRM、MCTS 等），一条是"agent 框架"（从 ReAct、Reflexion 等早期反应式规划，到多智能体协作、具身交互；基准如 Sokoban、FrozenLake、WebShop）。本文要把"基于 RL 的推理"和"结构化 agent 训练"在符号与语言任务上统一起来。

\textbf{结论与局限}：RAGEN 是一个在多轮、随机环境里用 RL 训练语言智能体的通用系统；它揭示了 agent 训练独有的新挑战——\textbf{梯度崩溃（gradient collapse）、rollout 漂移（rollout drift）、推理退化（reasoning degradation）}，并给出稳定训练的关键设计原则（rollout 过滤、梯度塑形、感知推理的奖励监督）。作者坦诚的\textbf{局限}：聚焦相对小规模任务、省略了 replay buffer 等成熟 RL 实践、未涉及多模态任务——留作未来工作。

%==================================================================
\section{核心 takeaways：给 AI 应用开发者的启示}
%==================================================================

\begin{enumerate}[leftmargin=1.6em,itemsep=4pt]
\item \textbf{多轮 agent RL 不是把单轮方法"加几轮"那么简单}。直接把 PPO/GRPO 搬过来，多半会\textbf{先涨后崩}。如果你打算自己训 agent，要预期到不稳定，并\textbf{提前准备稳定化手段}，而不是等崩了再救。

\item \textbf{把"奖励标准差、熵、梯度范数"当训练体检仪}。它们能在成功率掉之前给出预警：\emph{奖励 std / 熵异常波动}是早期信号，\emph{梯度范数尖峰}是"不可逆崩溃"的红线。在你的训练脚本里把这几个指标打到监控面板上，比只盯着 loss/成功率有用得多。

\item \textbf{"难样本优先"是个高性价比的稳定技巧}。StarPO-S 的核心——只在"模型时对时错（高奖励方差）"的实例上重点训练——既稳又省。这对应主动学习直觉，且实现简单（按 rollout 奖励 std 排序、保留 top $p\%$）。

\item \textbf{rollout 的"配方"很关键，有明确的调参甜区}：初始状态要\textbf{多样}、但每个 prompt 仍要有\textbf{多条 rollout 做对比}（甜区约 4 条/prompt）；每轮动作预算\textbf{适中}（5--6 个）；rollout 要\textbf{新鲜}（尽量在线、Online-1）。这些都是可以直接照搬的工程默认值。

\item \textbf{"让模型加 \texttt{<think>} 标签"并不等于"模型真的会推理"}。如果奖励只看最终成败、不看推理质量，模型会\textbf{偷懒退化}（推理越来越短）甚至\textbf{编幻觉推理却照样拿高分}。想让推理真正发挥作用，必须设计\textbf{感知推理、细粒度}的奖励（哪怕只是格式惩罚这种轻量手段，也比什么都不做强）。

\item \textbf{方法论价值大于"刷榜"价值}。这篇刻意用 0.5B 小模型 + 玩具环境，不是为了 SOTA，而是为了\textbf{把训练动力学看清楚}。对开发者的启示是：在投入大模型大算力之前，先在\textbf{极简可控环境}里验证你的训练配方是否稳定——这能省下大量试错成本。
\end{enumerate}

%==================================================================
\section*{参考文献（仅列最关键几条，英文原样）}
%==================================================================
\addcontentsline{toc}{section}{参考文献}
\begin{enumerate}[leftmargin=2em,itemsep=1pt]
\item J. Schulman, F. Wolski, P. Dhariwal, A. Radford, and O. Klimov. Proximal policy optimization algorithms, 2017. \url{https://arxiv.org/abs/1707.06347}.
\item Z. Shao, P. Wang, Q. Zhu, et al. DeepSeekMath: Pushing the limits of mathematical reasoning in open language models, 2024. \url{https://arxiv.org/abs/2402.03300}. （GRPO 出处）
\item DeepSeek-AI, D. Guo, et al. DeepSeek-R1: Incentivizing reasoning capability in LLMs via reinforcement learning, 2025. \url{https://arxiv.org/abs/2501.12948}.
\item Q. Yu, Z. Zhang, R. Zhu, et al. DAPO: An open-source LLM reinforcement learning system at scale, 2025. \url{https://arxiv.org/abs/2503.14476}. （KL 移除 + Clip-Higher 出处）
\item I. Shumailov, Z. Shumaylov, Y. Zhao, et al. AI models collapse when trained on recursively generated data. Nature, 631(8022):755--759, 2024. \url{https://doi.org/10.1038/s41586-024-07566-y}. （与 Echo Trap 相关）
\item S. Yao, H. Chen, J. Yang, and K. Narasimhan. WebShop: Towards scalable real-world web interaction with grounded language agents. NeurIPS, 2022.
\item Y. Zhou, A. Zanette, J. Pan, S. Levine, and A. Kumar. ArCHer: Training language model agents via hierarchical multi-turn RL, 2024. \url{https://arxiv.org/abs/2402.19446}.
\end{enumerate}

\end{document}
